\documentclass[aspectratio=169,11pt]{beamer}

\usetheme{Madrid}
\usecolortheme{default}
\usefonttheme{professionalfonts}
\setbeamertemplate{navigation symbols}{}
\setbeamertemplate{footline}[frame number]

\usepackage{amsmath, amssymb, booktabs}

\title{Synthesis, Frontier Questions, And Student Research Designs}
\subtitle{Turning Technology Questions Into Labor Research Designs}
\author{Special Topic 5: Technology, Innovation, and Labor -- Week 6}
\date{}

\begin{document}

\begin{frame}
  \titlepage
\end{frame}

\begin{frame}{Central question}
\begin{itemize}
    \item What makes a technology project a labor-economics paper?
    \item Start with the labor object: wages, employment, task content, mobility, training, job quality, bargaining power, or welfare.
    \item Then name the technology margin, adjustment channel, comparison, unit, timing, and incidence object.
    \item Week 6 goal: leave with project designs, not a generic future-of-work recap.
\end{itemize}
\end{frame}

\begin{frame}{Course synthesis}
\small
\begin{tabular}{p{0.24\linewidth} p{0.60\linewidth}}
\toprule
Week & Design object \\
\midrule
1. Tasks and frameworks & map technologies into tasks, skills, and equilibrium adjustment \\
2. Demand & distinguish displacement, augmentation, scale, and new tasks \\
3. Workers & study reskilling, mobility, career risk, and welfare costs \\
4. Firms & separate adoption, use, hiring, monitoring, and organization \\
5. Markets & trace rents, concentration, institutions, and distribution \\
6. Capstone & convert the sequence into credible research designs \\
\bottomrule
\end{tabular}
\end{frame}

\begin{frame}{Technology-to-labor research architecture}
\small
\begin{tabular}{p{0.30\linewidth} p{0.53\linewidth}}
\toprule
Design object & Required choice \\
\midrule
Technology margin & invention, adoption, use, exposure, diffusion, implementation \\
Labor outcome & wages, employment, vacancies, training, mobility, job quality, welfare \\
Worker channel & retraining, task switching, occupation moves, separations, retirement \\
Firm channel & substitution, augmentation, scale, hiring, supervision, job redesign \\
Market incidence & local markets, sectors, platforms, concentration, institutions \\
Welfare object & rents, risk, autonomy, bargaining, adjustment costs, distribution \\
\bottomrule
\end{tabular}
\vspace{0.25cm}
\begin{itemize}
    \item A paper that names a technology has not yet named a labor mechanism.
\end{itemize}
\end{frame}

\begin{frame}{Where the literature is}
\begin{itemize}
    \item Task frameworks and automation theory are mature: tasks, substitution, augmentation, and new work.
    \item Robot and local-exposure designs are the benchmark for market-level incidence.
    \item Worker-adjustment work measures earnings, mobility, training, and career risk better than subjective welfare or information frictions.
    \item Firm-adoption work now separates investment, adoption, use, skill demand, organization, and attitudes.
    \item Inequality work connects technology to rents, labor share, concentration, institutions, and worker voice.
\end{itemize}
\end{frame}

\begin{frame}{What remains poorly measured}
\begin{itemize}
    \item Actual worker use rather than access, exposure, or adoption announcements.
    \item Organizational implementation after adoption.
    \item Incumbent training versus firms hiring around incumbents.
    \item Algorithmic management, monitoring intensity, autonomy, and voice.
    \item Job quality, schedule risk, promotion ladders, and bargaining power.
    \item Compute, electricity, and data-center constraints as labor-market forces.
\end{itemize}
\end{frame}

\begin{frame}{Frontier directions}
\small
\begin{tabular}{p{0.30\linewidth} p{0.50\linewidth}}
\toprule
Frontier & Labor question \\
\midrule
Generative AI use & who gets augmented, displaced, trained, or monitored? \\
Firm reorganization & how do vacancies, teams, supervision, and promotion change? \\
Compute infrastructure & which places and workers gain from AI buildout and bottlenecks? \\
Global service trade & which tasks become tradable, automated, or newly valuable? \\
Measurement & which technology proxy predicts which labor outcome? \\
Worker voice & how do platform rules and algorithmic governance change bargaining? \\
\bottomrule
\end{tabular}
\end{frame}

\begin{frame}{Research design choices}
\begin{itemize}
    \item Labor object: what outcome matters, and why is it a labor outcome?
    \item Technology margin: invention, availability, adoption, use, exposure, or implementation?
    \item Unit: worker, occupation, firm, establishment, team, platform, place, or country?
    \item Counterfactual: less exposed units, later adopters, randomized rollout, policy shift, or pre-trends?
    \item Timing: short-run productivity, medium-run hiring, long-run mobility, and equilibrium incidence can differ.
    \item Welfare: wages and employment may miss risk, autonomy, voice, rents, and mobility costs.
\end{itemize}
\end{frame}

\begin{frame}{Common failure modes}
\begin{itemize}
    \item Technology-first framing with no labor mechanism.
    \item Treating a proxy as the technology itself.
    \item Confusing exposure with adoption or use.
    \item Ignoring selected adoption by firms, managers, or workers.
    \item Reading equilibrium incidence from partial-equilibrium evidence.
    \item Compressing timing into one coefficient.
    \item Interpreting wage effects as complete welfare effects.
\end{itemize}
\end{frame}

\begin{frame}{How to find a student project}
\begin{itemize}
    \item Avoid saturated broad questions; enter through a sharper mechanism.
    \item Replace "Does AI matter?" with "Which labor outcome, for which unit, through which adjustment channel?"
    \item Look for measurement gaps: use versus exposure, organization versus adoption, worker voice versus productivity.
    \item Use setting-specific evidence only when it reveals a portable labor mechanism.
    \item Make the contribution legible to labor economists who do not study the specific technology.
\end{itemize}
\end{frame}

\begin{frame}{Research Lab design}
\begin{itemize}
    \item Reproduce: build a synthetic local technology-exposure design inspired by robot-incidence research.
    \item Diagnose: classify what the measure captures and what it misses: adoption, use, exposure, timing, unit, welfare.
    \item Transfer: redesign the logic for generative AI worker use and firm adoption.
    \item Output: local exposure table, reduced-form incidence table, design diagnosis, and frontier transfer memo tables.
    \item Goal: practice turning technology measures into credible labor questions.
\end{itemize}
\end{frame}

\begin{frame}{Takeaways}
\begin{itemize}
    \item Technology-and-labor research begins with work, workers, firms, markets, institutions, and welfare.
    \item The technology measure must match the mechanism it claims to identify.
    \item Strong designs state unit, timing, comparison, adjustment channels, and welfare limits.
    \item Frontier opportunities are practical: better measurement, worker use, firm organization, infrastructure, global services, and voice.
    \item The course ends when students can turn a fast-moving topic into a researchable labor question.
\end{itemize}
\end{frame}

\end{document}
